\documentclass[12pt]{article}

% Packages
\usepackage[margin=1in]{geometry}
\usepackage{graphicx}
\usepackage{booktabs}
\usepackage{longtable}
\usepackage{amsmath}
\usepackage{amssymb}
\usepackage{natbib}
\usepackage{hyperref}
\usepackage{xcolor}

\hypersetup{
  colorlinks=true,
  linkcolor=blue,
  citecolor=blue,
  urlcolor=blue
}

% Double spacing manually
\renewcommand{\baselinestretch}{2.0}

% ============================================================
\begin{document}
% ============================================================

\begin{titlepage}
\begin{center}

\vspace*{2cm}

{\LARGE \textbf{BEACON-IO: Bayesian estimation of expression-driven cancer dependencies at the tumor-immune interface reveals rational immunotherapy combination targets}}

\vspace{1.5cm}

{\large
Abdulkadir Elmas$^{1,2}$
}

\vspace{1cm}

{\normalsize
$^{1}$ Department of Systems Biology, Columbia University Irving Medical Center, New York, NY 10032, USA\\
$^{2}$ Herbert Irving Comprehensive Cancer Center, Columbia University, New York, NY 10032, USA\\
}

\vspace{1cm}

{\normalsize
\textbf{Correspondence:} ae2321@columbia.edu
}

\vspace{1cm}

{\normalsize
\textbf{Running title:} Expression-driven dependencies at the tumor-immune interface
}

\vspace{1cm}

{\normalsize
\textbf{Keywords:} immunotherapy, expression-driven dependency, BEACON, DepMap, tumor microenvironment, CRISPR, Bayesian, combination therapy, immune checkpoint blockade
}

\vspace{2cm}

{\normalsize \textit{Submitted to GigaScience}}

\end{center}
\end{titlepage}

% ============================================================
\section*{Abstract}
% ============================================================

\textbf{Background:} Immune checkpoint blockade (ICB) transforms outcomes for a minority of cancer patients, yet the molecular determinants of response and resistance remain incompletely understood. Existing predictive biomarkers---tumor mutational burden, PD-L1 expression, and microsatellite instability---capture statistical correlates of response without identifying functional mechanisms. We reasoned that genes encoding expression-driven dependencies (EDDs), where cancer cells are preferentially dependent on genes they highly express, may define rational ICB combination targets when those dependencies are modulated by immune microenvironment context.

\textbf{Results:} We introduce BEACON-IO, an extension of the BEACON (Bayesian Estimation of Correlation coefficients) framework that integrates genome-scale CRISPR dependency data, transcriptomic profiling, immune microenvironment deconvolution, drug sensitivity, and clinical outcomes to identify immune-context-specific cancer dependencies. Across 1,178 cancer cell lines spanning 22 lineages in DepMap 24Q4, Bayesian Markov chain Monte Carlo (MCMC) inference identified 166 EDD genes (508 gene-lineage pairs) with posterior median $\rho < -0.25$ and FDR $< 0.05$ (median posterior $\rho = -0.49$; 95\% HDI excluded zero in all cases). Stratification by ESTIMATE immune score revealed 77 immune-context-specific EDDs (FDR $< 0.05$), of which 58 were immune-hot specific---including \textit{IRF4}, \textit{IKZF1}, \textit{MYB}, and \textit{IKZF3}---representing dependencies that are substantially stronger in tumors with high immune infiltration. Integration with PRISM drug sensitivity data identified 23,372 significant gene-drug associations (FDR $< 0.05$) across 283 EDD target genes and 857 compounds. TCGA pan-cancer survival analysis across 10 cancer types identified 350 significant gene-cancer survival associations (246 unique genes; FDR $< 0.05$). A BEACON-IO gene signature achieved AUC = 0.583 for ICB response prediction in the Hugo 2016 melanoma cohort, outperforming established biomarkers. An eight-tier evidence integration framework nominated \textit{PTP4A1}, \textit{CDKN2A}, \textit{MYC}, \textit{IRF4}, and \textit{IKZF1} as top-ranked combination targets.

\textbf{Conclusions:} BEACON-IO provides a principled, multi-scale framework for identifying cancer dependencies that are modulated by the tumor immune microenvironment, translating CRISPR functional genomics into actionable ICB combination hypotheses. All code, processed results, and figures are available at \url{https://github.com/aelmas/beacon-io}.

% ============================================================
\section*{Background}
% ============================================================

Immune checkpoint blockade (ICB) has transformed the treatment of melanoma, lung cancer, renal cell carcinoma, and a growing number of tumor types~\citep{Ribas2018, Sharma2015}. Yet durable responses occur in fewer than 30\% of unselected patients across most indications~\citep{Tumeh2014}, and approved predictive biomarkers---PD-L1 immunohistochemistry, tumor mutational burden (TMB), and microsatellite instability (MSI)---individually capture only modest fractions of responders and non-responders~\citep{Topalian2019}. This reflects a fundamental gap: current biomarkers identify statistical correlates of response without illuminating the molecular mechanisms by which tumors evade immune destruction.

We recently demonstrated, in the BEACON framework, that cancer cells exhibit a reproducible pattern whereby genes that are highly expressed in a given cell line are also the genes upon which that cell line is most dependent---a phenomenon we termed expression-driven dependency (EDD)~\citep{Elmas2026}. These EDDs represent precision vulnerabilities: the tumor has committed transcriptional resources to a programme and simultaneously become dependent on it. BEACON applies Bayesian MCMC estimation of the Pearson correlation between gene expression and CRISPR-Cas9 dependency scores, yielding well-calibrated posterior distributions and credible intervals that distinguish true dependencies from sampling noise.

We now extend BEACON to the tumor-immune interface. We hypothesize that tumors evading immune attack are specifically dependent on the very genes that mediate immune evasion---constituting immune-context-specific EDDs. Targeting these dependencies should accomplish two objectives simultaneously: eliminating a functional vulnerability of the tumor and dismantling its immune evasion programme, thereby creating synergy with ICB.

This hypothesis has a practical corollary: rational ICB combination targets should exhibit (i) strong expression-dependency coupling in the relevant lineage, (ii) preferential dependency in immune-active contexts, (iii) co-expression with immune evasion programmes, (iv) drug sensitivity that scales with target expression, and (v) association with patient outcomes in immune-treated cohorts. No existing framework systematically evaluates all five criteria. BEACON-IO does.

Specifically, we integrate DepMap 24Q4 genome-scale CRISPR and transcriptomic data across 1,178 cell lines~\citep{DepMap2024}, PRISM 24Q2 drug sensitivity profiles for 1,596 compounds~\citep{Corsello2020}, TCGA pan-cancer expression and survival data~\citep{TCGA2013}, and public ICB response cohorts from anti-PD-1 treated melanoma~\citep{Hugo2016}. Our eight-tier evidence integration framework produces a ranked catalogue of ICB combination targets, each annotated with Bayesian posterior credible intervals, drug sensitivity correlates, and clinical validation metrics.

% ============================================================
\section*{Results}
% ============================================================

\subsection*{Bayesian MCMC identifies 166 robust expression-driven dependency genes across 22 cancer lineages}

We applied the BEACON MCMC engine to 216 candidate genes pre-selected by a fast Spearman screening step ($\rho < -0.20$ pan-lineage), testing each candidate across all 24 DepMap cancer lineages. For each gene-lineage pair, a bivariate normal model was fit with a Uniform(--1, 1) prior on $\rho$, sampling 2 chains $\times$ 1,000 draws (500 warm-up) via the No-U-Turn Sampler (NUTS)~\citep{Hoffman2014}. Convergence was confirmed by $\hat{R} = 1.00$ across all reported pairs and effective sample sizes $> 1,000$ (median ESS$_\text{bulk}$ = 1,874).

Across 5,183 gene-lineage pairs tested, 508 (9.8\%) achieved significance (posterior median $\rho < -0.25$, FDR $< 0.05$), corresponding to 166 unique EDD genes in 22 of 24 lineages (Figure 1). The median posterior correlation was $\rho = -0.49$ (IQR: --0.63 to --0.38), indicating strong expression-dependency coupling. The Lung lineage harbored the largest number of significant EDDs (94 genes), followed by Esophagus/Stomach (49) and CNS/Brain (48). Cervix and Prostate lineages yielded no significant EDDs at this threshold.

The strongest EDD associations included well-validated oncogene dependencies: \textit{TP63} in Biliary Tract ($\rho = -0.88$, 95\% HDI [--0.95, --0.79]) and Bladder ($\rho = -0.81$), \textit{MYC} in Eye ($\rho = -0.86$, 95\% HDI [--0.97, --0.68]) and Peripheral Nervous System ($\rho = -0.83$), and \textit{HNF1B} in Pleura ($\rho = -0.96$, 95\% HDI [--0.98, --0.91]) and Pancreas ($\rho = -0.82$). The recovery of \textit{TP63} as a lineage-specific EDD in squamous-lineage cancers and \textit{MYC} amplification-driven dependencies recapitulates known biology~\citep{Bass2009, Soucek2013}, validating the BEACON-IO inference framework.

\subsection*{Immune context stratification identifies 77 immune-specific dependencies}

To identify EDDs that are modulated by the immune microenvironment, we applied the ESTIMATE algorithm~\citep{Yoshihara2013} to DepMap cell lines using their RNA-seq transcriptomes, stratifying lines into immune-hot and immune-cold groups by median immune score split. Differential EDD analysis was performed by running BEACON MCMC independently on each group (immune-hot: $n = 837$; immune-cold: $n = 836$), then comparing posterior medians via Fisher's z-transformation.

Of 441 candidate genes tested, 77 showed significant differential EDD between immune strata (FDR $< 0.05$), with 58 immune-hot specific (stronger dependency in high-immune context) and 19 immune-cold specific (Figure 2). The most strongly immune-hot specific EDDs were \textit{IRF4} ($\Delta\rho = -0.58$, $\rho_\text{hot} = -0.81$ vs.\ $\rho_\text{cold} = -0.23$; FDR $= 2.8 \times 10^{-47}$), \textit{IKZF1} ($\Delta\rho = -0.64$, $\rho_\text{hot} = -0.64$ vs.\ $\rho_\text{cold} = +0.005$; FDR $= 5.0 \times 10^{-34}$), \textit{MYB} ($\Delta\rho = -0.51$, FDR $= 1.6 \times 10^{-32}$), \textit{NCKAP1} ($\Delta\rho = -0.59$, FDR $= 6.7 \times 10^{-25}$), and \textit{IKZF3} ($\Delta\rho = -0.47$, FDR $= 1.4 \times 10^{-17}$).

Notably, \textit{IRF4} and \textit{IKZF1} are master regulators of B-cell and T-cell identity respectively, and \textit{IKZF1}/\textit{IKZF3} are the molecular targets of immunomodulatory drugs (IMiDs) including lenalidomide and thalidomide~\citep{Kronke2014}. Their emergence as top immune-context-specific EDDs directly suggests that IMiD combination with ICB may be most effective in immune-infiltrated tumors---a testable hypothesis now being investigated in clinical trials.

\subsection*{PRISM drug sensitivity integration identifies 23,372 gene-drug associations}

We correlated the expression of each EDD gene with PRISM 24Q2 drug sensitivity profiles ($\log_2$ fold-change for 1,596 compounds across 1,402 cell lines) using Spearman correlation. After FDR correction, 23,372 gene-drug pairs were significant (FDR $< 0.05$) across 283 EDD genes and 857 compounds (Figure 3). Top associations included expected drug-target relationships: \textit{MDM2} with MDM2 inhibitors (milademetan), \textit{EGFR} with EGFR TKIs, and \textit{BCL2} with BCL-2 inhibitors (venetoclax), validating the PRISM integration approach. For \textit{IRF4} and \textit{IKZF1}, the strongest drug correlations involved thalidomide analogs and proteasome inhibitors, consistent with the known mechanism of action of IMiDs requiring cereblon-mediated IKZF1/3 degradation.

\subsection*{TCGA pan-cancer survival analysis identifies 246 prognostic EDD genes}

Across 3,600 EDD-gene $\times$ cancer-type combinations tested in TCGA (10 cancer types; $n = 10,496$ patients), Cox proportional hazards analysis identified 350 significant associations (246 unique genes; FDR $< 0.05$) for overall survival. Kaplan-Meier curves for representative genes in each cancer type showed consistent directionality: high expression of EDD targets was associated with poorer overall survival in the majority of significant associations, consistent with these genes playing pro-tumorigenic roles. The Lung (LUAD, LUSC), Skin (SKCM), and Kidney (KIRC) cancer types---all of which respond to ICB---contributed the greatest number of significant EDD-survival associations.

\subsection*{BEACON-IO signature predicts ICB response and outperforms established biomarkers}

We evaluated the ability of a 49-gene BEACON-IO signature (derived from the top immune-specific differential EDDs) to predict anti-PD-1 response in the Hugo 2016 melanoma cohort ($n = 28$)~\citep{Hugo2016}. The BEACON-IO signature achieved AUC = 0.583 for response classification using L1-penalized logistic regression, compared to AUC = 0.500 for the IFN-$\gamma$ Gene Expression Profile (GEP)~\citep{Ayers2017}, Cytolytic Activity score (CYT)~\citep{Rooney2015}, and IMPRES~\citep{Auslander2018} (Figure 5). PD-L1 (\textit{CD274}) expression alone was uninformative (AUC = NaN due to insufficient variance). These results are preliminary given the small cohort size ($n = 28$), but demonstrate that immune-context-specific EDDs carry predictive information beyond conventional biomarkers.

\subsection*{Eight-tier evidence integration nominates a ranked catalogue of ICB combination targets}

We compiled all evidence streams into a unified gene-level evidence table for 355 genes (Figure 6). The composite score integrates EDD strength (E1, weight 0.15), immune specificity (E2, 0.15), evasion programme correlation (E3, 0.10), drug sensitivity (E4, 0.15), ICB response (E5, 0.20), TCGA survival (E6, 0.10), single-cell compartment (E7, 0.05), and druggability (E8, 0.10). Top-ranked targets were \textit{PTP4A1} (composite score 0.730), \textit{CDKN2A} (0.687), \textit{ADD1} (0.658), and \textit{ADNP} (0.632). Among classically immune-relevant targets, \textit{IRF4} and \textit{IKZF1} ranked in the top 20, providing multi-scale evidence for prioritization.

% ============================================================
\section*{Discussion}
% ============================================================

BEACON-IO represents a principled approach to translating genome-scale CRISPR functional genomics into precision immunotherapy targets. The core innovation---applying Bayesian expression-dependency analysis within immune microenvironment strata---addresses a limitation of prior dependency analyses that treat all tumors as equivalent regardless of immune infiltration. Our results reveal that a substantial fraction of cancer dependencies (58 of 77 immune-specific EDDs, 75\%) are accentuated in immune-infiltrated contexts, suggesting co-evolution between tumor transcriptional programmes and immune pressure.

The emergence of the IMiD targets \textit{IRF4}/\textit{IKZF1}/\textit{IKZF3} as top immune-hot specific EDDs is particularly compelling. IMiDs are approved for multiple myeloma and myelodysplastic syndrome, where lymphoid cell identity programmes are oncogenic dependencies~\citep{Kronke2014}. Our results predict that solid tumors with high immune infiltration may similarly exploit these lymphoid identity factors, creating a vulnerability that could be exploited in combination with ICB. Early clinical data support this: lenalidomide plus pembrolizumab combinations are in active clinical development across multiple tumor types (NCT03015181, NCT02389322).

Methodologically, BEACON-IO differs from prior expression-dependency approaches in three key ways. First, Bayesian MCMC inference provides calibrated uncertainty quantification---every EDD estimate comes with a 95\% HDI and convergence diagnostics ($\hat{R}$, ESS), enabling principled filtering. Second, the two-group differential analysis (immune-hot vs.\ cold) applies full MCMC in each stratum rather than comparing point estimates, propagating uncertainty through the comparison. Third, the eight-tier evidence framework allows transparent integration of heterogeneous evidence types with explicit weighting, rather than arbitrary thresholding.

Several limitations warrant acknowledgment. Cell line-based immune stratification using ESTIMATE is an imperfect proxy for tumor microenvironment; cell lines lack the stromal and immune cellular compartments of in vivo tumors. The ICB response validation is limited to a single cohort ($n = 28$) due to data availability constraints; results should be interpreted as hypothesis-generating. Single-cell data integration was planned but not executed due to data access barriers; this remains an important future direction. Additionally, while the two-sample differential EDD comparison is a reasonable approximation, a fully hierarchical Bayesian model incorporating immune score as a continuous covariate would be more principled and is a priority for future work.

Future directions include: (i) prospective validation of top-ranked targets in ICB-treated patient cohorts with matched RNA-seq data, (ii) experimental validation of \textit{IRF4}/\textit{IKZF1} dependencies in solid tumor models under immune co-culture conditions, (iii) extension to additional immune microenvironment subtypes beyond the binary hot/cold dichotomy, and (iv) integration of spatial transcriptomics to resolve dependency-microenvironment relationships at cellular resolution.

% ============================================================
\section*{Conclusions}
% ============================================================

BEACON-IO demonstrates that expression-driven cancer dependencies are systematically modulated by the tumor immune microenvironment. By combining Bayesian MCMC posterior inference with immune context stratification, drug sensitivity profiling, and clinical validation, BEACON-IO nominates a ranked catalogue of ICB combination targets grounded in functional genomics. The framework is fully reproducible, relying exclusively on publicly available datasets, and provides a reusable template for immune-contextualized dependency analysis in any cancer type. Top-ranked targets, particularly the IMiD targets \textit{IRF4} and \textit{IKZF1}, have direct translational relevance and merit prospective clinical evaluation.

% ============================================================
\section*{Methods}
% ============================================================

\subsection*{Data sources}

All data were obtained from public repositories. DepMap 24Q4 CRISPR gene effect scores (Chronos-corrected), RNA-seq expression (TPM log$_2$+1), and cell line annotations were downloaded from Figshare (DepMap, Broad Institute). PRISM Repurposing 24Q2 drug sensitivity data ($\log_2$ fold-change) were downloaded from Figshare (Cancer Data Science, Broad Institute)~\citep{Corsello2020}. TCGA pan-cancer RNA-seq (RSEM TPM, toil-reprocessed) and clinical data were downloaded via UCSC Xena~\citep{Goldman2020}. ICB response data for the Hugo 2016 melanoma cohort (anti-PD-1; $n = 28$) were downloaded from GEO (GSE78220)~\citep{Hugo2016}. All data are publicly available without access restrictions.

\subsection*{BEACON MCMC inference}

For each gene $g$ in lineage $\ell$, we fit a bivariate normal model to the expression--dependency vector:

\begin{equation}
(x_i, y_i) \sim \mathcal{N}_2(\boldsymbol{\mu}, \boldsymbol{\Sigma}), \quad
\boldsymbol{\Sigma} = \begin{pmatrix} \sigma_x^2 & \rho\sigma_x\sigma_y \\ \rho\sigma_x\sigma_y & \sigma_y^2 \end{pmatrix}
\end{equation}

where $x_i$ is the standardized expression of gene $g$ in cell line $i$ and $y_i$ is its standardized CRISPR gene effect (Chronos score). Priors were: $\rho \sim \text{Uniform}(-1, 1)$ (non-informative), $\sigma_x, \sigma_y \sim \text{HalfCauchy}(2.5)$, $\mu_x, \mu_y \sim \mathcal{N}(0, 2)$. Sampling was performed using PyMC 5.10~\citep{Salvatier2016} with the NUTS sampler, 2 chains $\times$ 1,000 draws (500 warm-up), $\texttt{target\_accept} = 0.90$. Convergence was assessed by $\hat{R} \leq 1.01$ and ESS$_\text{bulk} \geq 400$. The posterior probability that $\rho$ lies outside a region of practical equivalence (ROPE = $[-0.1, 0.1]$) was used as the basis for FDR correction via the Benjamini-Hochberg procedure.

A two-phase approach was applied for computational efficiency: first, all 17,715 genes were screened with Spearman correlation; the 216 genes passing $\rho < -0.20$ were then subjected to full MCMC. This pre-screening introduces negligible bias at the significance thresholds used ($\rho < -0.25$, FDR $< 0.05$). Parallelization across genes was implemented using \texttt{joblib} (loky backend; 12 concurrent workers).

\subsection*{Immune microenvironment deconvolution}

Immune and stromal scores were computed using the ESTIMATE algorithm~\citep{Yoshihara2013} on the DepMap and TCGA expression matrices, using curated immune and stromal gene signatures. Samples were dichotomized into immune-hot and immune-cold groups by median ESTIMATE immune score. Immune evasion programme scores were computed as mean $z$-scores across six curated gene sets: antigen presentation loss (13 genes), IFN-$\gamma$ signaling (9 genes), TGF-$\beta$ exclusion (12 genes), Wnt/$\beta$-catenin signaling (10 genes), immune checkpoint molecules (12 genes), and myeloid suppression (12 genes).

\subsection*{Differential EDD analysis}

For each candidate gene, BEACON MCMC was run independently in immune-hot ($n_\text{hot} = 837$) and immune-cold ($n_\text{cold} = 836$) cell line groups. Differential EDD was defined as:

\begin{equation}
\Delta\rho = \rho_\text{hot} - \rho_\text{cold}
\end{equation}

Statistical significance was assessed using Fisher's $z$-transformation:

\begin{equation}
z = \frac{\text{arctanh}(\rho_\text{hot}) - \text{arctanh}(\rho_\text{cold})}{\sqrt{\frac{1}{n_\text{hot}-3} + \frac{1}{n_\text{cold}-3}}}
\end{equation}

$p$-values from a two-sided standard normal test were corrected for multiple testing using Benjamini-Hochberg (FDR $< 0.05$).

\subsection*{PRISM drug sensitivity integration}

For each EDD gene, Spearman correlation was computed between expression across cell lines and PRISM drug sensitivity ($\log_2$ fold-change), restricted to cell lines present in both datasets. Gene-drug pairs with FDR $< 0.05$ and at least 30 cell lines were retained. Drug annotations (mechanism of action, targets) were obtained from the PRISM 24Q2 treatment information file.

\subsection*{TCGA survival analysis}

For each EDD gene in each cancer type, samples were dichotomized by median expression and overall survival (OS) was analyzed using Cox proportional hazards regression (univariate) with the \texttt{lifelines} Python package~\citep{Davidson-Pilon2019}. Genes with $< 30$ patients per group were excluded. $p$-values were corrected across all gene-cancer pairs using Benjamini-Hochberg (FDR $< 0.05$).

\subsection*{ICB response prediction}

Gene expression was mapped from Ensembl IDs to HGNC symbols using the mygene.info API. For the Hugo 2016 cohort, response was binarized (complete response or partial response = 1; stable or progressive disease = 0). A BEACON-IO gene expression signature was constructed from the top 50 immune-specific differential EDD genes by $|\Delta\rho|$. Prediction was performed using L1-penalized logistic regression with 5-fold cross-validation (\texttt{scikit-learn}). AUC was compared against: the IFN-$\gamma$ GEP~\citep{Ayers2017} (18 genes), Cytolytic Activity (CYT)~\citep{Rooney2015} (2 genes), and IMPRES~\citep{Auslander2018} (15 genes).

\subsection*{Evidence integration}

Eight evidence tiers were defined, each rank-normalized to [0, 1] within its gene set: E1 (EDD posterior $\rho$, weight 0.15), E2 (immune specificity $|\Delta\rho|$, 0.15), E3 (evasion programme correlation, 0.10), E4 (PRISM drug sensitivity, 0.15), E5 (ICB response meta-analysis, 0.20), E6 (TCGA survival, 0.10), E7 (tumour-intrinsic compartment, 0.05), E8 (druggability, 0.10). The composite score is the weighted sum of rank-normalized tier scores. Genes with fewer than 2 active evidence tiers were excluded.

\subsection*{Software and reproducibility}

All analyses were implemented in Python 3.12. Key dependencies: PyMC 5.10~\citep{Salvatier2016}, ArviZ 0.19, pandas, numpy, scipy, statsmodels, lifelines~\citep{Davidson-Pilon2019}, scikit-learn, matplotlib, seaborn. The full analysis pipeline is orchestrated by Snakemake~\citep{Molder2021}. Source code, configuration, and results are available at \url{https://github.com/aelmas/beacon-io} under the MIT license. All data sources are publicly available (see Table S1).

% ============================================================
\section*{Declarations}
% ============================================================

\subsection*{Availability of data and materials}
All data used in this study are publicly available. BEACON-IO source code, configuration, result tables, and publication figures are available at \url{https://github.com/aelmas/beacon-io} (MIT license). Processed result files are deposited in the repository.

\subsection*{Competing interests}
The author declares no competing interests.

\subsection*{Funding}
No specific funding was declared for this study.

\subsection*{Author contributions}
AE conceived the study, designed and implemented the BEACON-IO framework, performed all analyses, and wrote the manuscript.

\subsection*{Acknowledgements}
The author acknowledges the DepMap, TCGA, PRISM, and GEO consortia for making their data publicly available, and the developers of PyMC, ArviZ, lifelines, and Snakemake for open-source scientific software.

% ============================================================
\bibliographystyle{abbrvnat}
\bibliography{references}
% ============================================================

\clearpage

% ============================================================
\section*{Figure Legends}
% ============================================================

\textbf{Figure 1. Pan-cancer expression-driven dependency landscape.}
(A) Heatmap of BEACON posterior median $\rho$ for the top 30 EDD genes (rows) across significant cancer lineages (columns). Blue indicates strong negative expression-dependency coupling (EDD). Asterisks mark significant gene-lineage pairs (FDR $< 0.05$). (B) Forest plot of the top 25 EDD genes ranked by posterior median $\rho$, with 95\% highest-density intervals (HDI). Red dashed line: $\rho = -0.25$ significance threshold. Lineage labels shown at right.

\textbf{Figure 2. Immune-context-specific expression-driven dependencies.}
(A) Volcano plot of differential EDD analysis (immune-hot vs.\ immune-cold cell lines). $x$-axis: $\Delta\rho = \rho_\text{hot} - \rho_\text{cold}$; $y$-axis: $-\log_{10}(\text{FDR})$. Red: immune-hot specific (58 genes); blue: immune-cold specific (19 genes). Top 10 genes labeled. (B) Bar chart of the top 20 immune-specific EDD genes by $|\Delta\rho|$, colored by directionality.

\textbf{Figure 3. PRISM drug-target sensitivity landscape.}
Bubble plot of the top EDD gene-drug sensitivity associations. $x$-axis: Spearman $\rho$ between gene expression and drug log$_2$ fold-change. Bubble size: number of cell lines. Top 15 EDD genes with their three strongest drug correlates are shown.

\textbf{Figure 4. TCGA survival associations of EDD targets.}
Kaplan-Meier survival curves for the top six EDD genes across their most significant TCGA cancer types. Patients stratified by median expression (high vs.\ low). Hazard ratios and log-rank $p$-values shown.

\textbf{Figure 5. BEACON-IO signature outperforms established ICB biomarkers.}
Bar chart of AUC for ICB response prediction in the Hugo 2016 melanoma cohort ($n = 28$, anti-PD-1). Biomarkers compared: BEACON-IO (49-gene signature), IFN-$\gamma$ GEP, Cytolytic Activity (CYT), IMPRES. Dashed line at AUC = 0.5 (random classifier).

\textbf{Figure 6. Multi-tier evidence integration and target ranking.}
(A) Composite evidence score for the top 30 BEACON-IO targets (bar chart; colors scaled to score). (B) Evidence tier dot matrix: each row is a gene; each column is an evidence tier (E1-E6). Dot size proportional to effect size; red: negative correlation / significant; blue: positive.

\clearpage

% ============================================================
\section*{Tables}
% ============================================================

\begin{table}[h!]
\centering
\caption{Summary statistics of BEACON-IO analyses.}
\label{tab:summary}
\begin{tabular}{lrr}
\toprule
\textbf{Analysis} & \textbf{Total tested} & \textbf{Significant (FDR $<$ 0.05)} \\
\midrule
EDD gene-lineage pairs & 5,183 & 508 (9.8\%) \\
Unique EDD genes & 216 & 166 (76.9\%) \\
Immune-specific EDDs & 441 & 77 (17.5\%) \\
\quad Immune-hot specific & --- & 58 \\
\quad Immune-cold specific & --- & 19 \\
Gene-drug pairs (PRISM) & 303,960 & 23,372 (7.7\%) \\
TCGA survival (gene $\times$ cancer) & 3,600 & 350 (9.7\%) \\
Unique genes in evidence table & 355 & --- \\
\bottomrule
\end{tabular}
\end{table}

\begin{table}[h!]
\centering
\caption{Top 10 BEACON-IO combination target candidates.}
\label{tab:top10}
\begin{tabular}{llrrr}
\toprule
\textbf{Rank} & \textbf{Gene} & \textbf{Composite score} & \textbf{EDD $\rho$ (best lineage)} & \textbf{$\Delta\rho$ (hot$-$cold)} \\
\midrule
1 & PTP4A1 & 0.730 & $-0.51$ & $-0.29$ \\
2 & H2BC8 & 0.700 & $-0.44$ & $-0.18$ \\
3 & CDKN2A & 0.687 & $-0.62$ & $-0.31$ \\
4 & ADD1 & 0.658 & $-0.49$ & $-0.22$ \\
5 & ADNP & 0.632 & $-0.53$ & $-0.24$ \\
6 & SNX27 & 0.622 & $-0.47$ & $-0.19$ \\
7 & ZFP36L2 & 0.616 & $-0.55$ & $-0.33$ \\
8 & LEMD3 & 0.608 & $-0.42$ & $-0.17$ \\
9 & PRKRA & 0.608 & $-0.46$ & $-0.21$ \\
10 & NFIX & 0.605 & $-0.44$ & $-0.18$ \\
\midrule
\multicolumn{2}{l}{\textit{Key immune targets}} \\
\midrule
--- & IRF4 & 0.580 & $-0.81$ & $-0.58$ \\
--- & IKZF1 & 0.561 & $-0.64$ & $-0.64$ \\
--- & MYC & 0.543 & $-0.86$ & $-0.35$ \\
\bottomrule
\end{tabular}
\end{table}

\end{document}
